\section{Resultados obtenidos}

Gracias al trabajo realizado durante el periodo de residencias profesionales en TenarisTamsa, en colaboración con el equipo de Operational Planning IT, se desarrolló un sistema de software orientado a la implementación de aplicaciones web mediante tecnologías modernas para el Sistema Integral de Producción.

La implementación de este sistema ha generado un impacto significativo en los procesos productivos ejecutados por los operadores durante las jornadas laborales, mejorando la eficiencia en la captura, consulta y administración de la información operativa.

De acuerdo con las políticas de privacidad y confidencialidad establecidas por la empresa, toda la información sensible presente en las imágenes y evidencias mostradas a continuación ha sido ocultada.

Con base en los resultados obtenidos tras la implementación piloto del sistema en dos plantas, se ha demostrado una reducción de hasta un 35\% en la pérdida de información, gracias a la digitalización de reportes y procesos. Esto representa una mejora considerable en comparación con el método anterior, en el cual los operadores realizaban el llenado manual de los formatos físicos.

Asimismo, se logró disminuir el consumo de papel en aproximadamente un 10\% dentro de las plantas participantes, reduciendo los costos operativos asociados y contribuyendo a la disminución de desperdicios generados por las distintas áreas.

Finalmente, el sistema permitió la unificación de la información entre ambas plantas donde se realizaron las pruebas piloto, facilitando el acceso a datos de manera más rápida, precisa y eficiente. Esto ha permitido compartir áreas de mejora y optimizar la toma de decisiones entre los responsables de cada planta, fortaleciendo la colaboración y estandarización de los procesos operativos.
